\begin{textsample}{POS Dim 5 – human – Score 1.00 – es/es\_ef\_anos\_iniciais\_ciencias\_humanas\_geo\_1\_p2.txt}  \label{ex:f5_pos_004}
Geografia, o que é?
Professores indígenas do Acre “ Geografia é onde o \textbf{rio} está.
Onde o município está. É para onde vem o sol. É para onde vai o sol.
Este \textbf{rio} para onde vai? É divisão das águas. É igarapé, igapó, lago, açude, mar. É a medição da terra, a demarcação. É fotografia, desenho, cor, é um mapa.
Geografia é o entendimento da aldeia e do mundo.
Do nosso mundo e do mundo do branco. É a cidade, o Brasil e os outros países. É a história do mundo.
O mundo é a terra, a terra é a aldeia.
O \textbf{rio} que cai num outro \textbf{rio}.
Que cai no mar.
Geografia é o depois do mar... ”

% matched lemmas: rio
\end{textsample}
